\documentclass[11pt,a4paper]{article}
\usepackage[margin=25mm]{geometry}
\usepackage{amsmath,amssymb,booktabs,siunitx,hyperref,listings,microtype}
\hypersetup{colorlinks=true,linkcolor=blue,urlcolor=blue,citecolor=blue,pdftitle={Report 4 - Turbine Model Family for OpenHPL/OpenHPLjl},pdfauthor={Madhusudhan Pandey}}
\lstdefinelanguage{Julia}{keywords={using,begin,end,function,return,module,export,include},sensitive=true,comment=[l]\#,morestring=[b]"}
\lstset{language=Julia,basicstyle=\ttfamily\small,breaklines=true,frame=single,columns=fullflexible,showstringspaces=false}
\title{Report 4 - Turbine Model Family for OpenHPL/OpenHPLjl}
\author{Madhusudhan Pandey}
\date{19 September 2026}
\begin{document}
\maketitle
\begin{abstract}
This report introduces the turbine model family in OpenHPLjl. The goal is to progress from transparent analytical hydraulic-to-mechanical conversion toward guide-vane, characteristic-map, and turbine-type-specific models while preserving one reusable hydraulic interface. The structure follows the same concise engineering sequence used for the reservoir, rigid-pipe, and surge-tank reports.
\end{abstract}

\section{Updated Context}
The turbine is the energy-conversion component between the hydraulic waterway and the rotating shaft:
\[
\boxed{\text{Reservoir}\rightarrow\text{RigidPipe}\rightarrow\text{SurgeTank}\rightarrow\text{Turbine}\rightarrow\text{Shaft}\rightarrow\text{Generator}}.
\]
The waterway determines head and flow. The turbine converts hydraulic power into mechanical power. OpenHPL includes turbine models and supporting turbine documentation, including Francis-turbine formulations and characteristic relationships \cite{Sharefi,OpenHPLTurbine}.

\section{Generalized Concept Model}
The generalized turbine receives hydraulic head and flow and produces mechanical power:
\[
\boxed{H,\;Q,\;y,\;\omega \rightarrow \text{Turbine} \rightarrow P_m,\;T_m}.
\]
The hydraulic head drop is
\[
H=H_{\rm in}-H_{\rm out}.
\]
The basic hydraulic power is
\[
P_h=\rho gQH.
\]
Mechanical power is
\[
\boxed{P_m=\eta \rho gQH}.
\]
When shaft speed is available,
\[
T_m=\frac{P_m}{\omega}.
\]

\section{Other Models Available in the Literature}
A useful hierarchy is
\[
\boxed{\text{ideal power converter}\rightarrow\text{constant-efficiency turbine}\rightarrow\text{gate-flow model}\rightarrow\text{lookup/Hill-chart model}\rightarrow\text{type-specific model}}.
\]
The simplest model computes power from network flow. A gate model adds guide-vane dependence. Higher-fidelity Francis, Pelton, and Kaplan models use turbine-specific nondimensional relationships and characteristic maps. OpenHPL documentation includes turbine characteristic relationships and Francis-turbine modeling resources \cite{Sharefi,OpenHPLTurbine}.

\section{Mathematics Involved in the OpenHPL Concept/Model}
For the first model family,
\[
H=H_{\rm in}-H_{\rm out},
\qquad
Q=Q_{\rm in},
\qquad
Q_{\rm out}=-Q.
\]
The ideal/constant-efficiency conversion is
\[
\boxed{P_m=\rho g\eta QH}.
\]
A first analytical guide-vane approximation is
\[
\boxed{Q=K_q y\sqrt{H}},
\]
for positive head. Combining the equations gives
\[
P_m=\rho g\eta K_q y H^{3/2}.
\]
In later models, $\eta$ and $Q$ become nonlinear functions of guide-vane opening, speed, and head:
\[
Q=f_Q(H,\omega,y),
\qquad
\eta=f_\eta(H,\omega,Q,y).
\]

\section{OpenHPL-Specific Modeling Interpretation}
The OpenHPL turbine philosophy is compatible with a progression from simple algebraic conversion toward nonlinear characteristic-map models. The turbine should remain connected hydraulically through the same head/flow ports while the internal constitutive equations become more detailed. This allows a simple turbine to be replaced by a Francis, Pelton, Kaplan, or lookup-table model without redesigning the surrounding waterway.

\section{Implementation in Julia ModelingToolkit / SciML}
ModelingToolkit supports hierarchical components and acausal connectors. The current package therefore places turbine alternatives under one model-family folder.

\begin{lstlisting}
src/Turbines/
|-- IdealTurbine.jl
|-- SimpleGateTurbine.jl
|-- TurbineLookup.jl      # planned
|-- FrancisTurbine.jl     # planned
|-- PeltonTurbine.jl      # planned
|-- KaplanTurbine.jl      # planned
\end{lstlisting}

The first conversion model is

\begin{lstlisting}
@component function IdealTurbine(;
    name, rho=1000.0, g=9.81, eta=1.0)

    @named inlet = HydraulicPort()
    @named outlet = HydraulicPort()

    @variables Q(t) H(t) Pm(t)

    eqs = [
        Q ~ inlet.Q
        outlet.Q ~ -Q
        H ~ inlet.H - outlet.H
        Pm ~ rho*g*eta*Q*H
    ]

    System(eqs, t, [Q,H,Pm], [rho,g,eta];
        systems=[inlet,outlet], name=name)
end
\end{lstlisting}

The first guide-vane model adds
\[
Q=K_q y\sqrt{H}.
\]

\section{Model Assembly and Connections}
A turbine can be connected downstream of the hydraulic system:
\begin{lstlisting}
@named tank = SurgeTank()
@named penstock = RigidPipe()
@named turbine = SimpleGateTurbine(y = 0.5)

connections = [
    connect(tank.outlet, penstock.inlet)
    connect(penstock.outlet, turbine.inlet)
]
\end{lstlisting}
A tail-water boundary is connected to the turbine outlet. A future rotational connector will connect turbine torque to the shaft.

\section{Numerical Experiment / Validation}
For
\[
H=100~\si{\meter},\qquad y=0.5,\qquad K_q=1,
\]
the simple gate law gives
\[
Q=1\cdot0.5\sqrt{100}=5~\si{\meter\cubed\per\second}.
\]
With
\[
\rho=1000~\si{\kilogram\per\meter\cubed},\qquad
g=9.81~\si{\meter\per\second\squared},\qquad
\eta=0.90,
\]
the expected mechanical power is
\[
\boxed{P_m=4.4145~\si{\mega\watt}}.
\]
This provides a simple deterministic unit-test target before introducing characteristic maps.

\section{Expected Outputs}
The minimum turbine outputs are
\[
Q(t),\qquad H(t),\qquad P_m(t),
\]
and later
\[
\eta(t),\qquad T_m(t),\qquad \omega(t),\qquad y(t).
\]
Useful plots include mechanical power versus guide-vane opening, efficiency versus operating point, head-flow response, and turbine characteristic maps.

\section{Engineering Interpretation}
The turbine family separates model fidelity from system architecture. The surrounding hydropower plant sees the same hydraulic interface while the internal model can evolve from
\[
\boxed{\text{ideal}\rightarrow\text{gate-controlled}\rightarrow\text{lookup map}\rightarrow\text{Francis/Pelton/Kaplan}}.
\]
This is important for FCR and FREKI studies because model complexity can be increased only where measurements and validation justify it.

\section*{Conclusion}
The first turbine family establishes the hydraulic-to-mechanical conversion
\[
\boxed{P_m=\rho g\eta QH}
\]
and a simple gate-flow relation
\[
\boxed{Q=K_qy\sqrt{H}}.
\]
The next turbine milestone is a lookup-table/Hill-chart model, followed by a rotational connector and shaft dynamics.

\begin{thebibliography}{9}
\bibitem{Sharefi}
OpenHPL resources, Francis turbine and hydraulic-system formulations.
\url{https://build.openmodelica.org/Documentation/OpenHPL%203.0.1/Resources/Documents/Sharefi2011.pdf}

\bibitem{OpenHPLTurbine}
OpenHPL turbine-model resource.
\url{https://build.openmodelica.org/Documentation/OpenHPL%203.0.1/Resources/Documents/Turbines_model.pdf}

\bibitem{MTK}
SciML, ModelingToolkit.jl acausal component modeling.
\url{https://docs.sciml.ai/ModelingToolkit/stable/tutorials/acausal_components/}
\end{thebibliography}
\end{document}
